\documentclass[11pt]{article}
\usepackage[a4paper,margin=23mm]{geometry}
\usepackage{amsmath,amssymb,booktabs,url}
\title{VNAV Lab 4: Trajectory Optimization\\Deliverables 1 and 2}
\author{Team Fullcans}\date{}
\begin{document}\maketitle
\section*{Deliverable 1: Single-segment minimum velocity}
Let $P(t)=\sum_{i=1}^N p_i t^i$, so $P(0)=0$ is already imposed.
Then
\[
J=\int_0^1\!\left(\sum_{i=1}^N i p_i t^{i-1}\right)^2dt
=\boldsymbol p^T Q\boldsymbol p,\qquad Q_{ij}=\frac{ij}{i+j-1},\quad i,j=1,\ldots,N.
\]
The endpoint position is $A\boldsymbol p=b$, with $A=[1,\ldots,1]$, $b=1$.
Here the objective is $p^TQp$; a solver using $\tfrac12p^THp$ needs $H=2Q$.
\subsection*{1. Degree one}
$P(1)=p_1=1$, hence $P^*(t)=t$ and $J^*=1$.
\subsection*{2. Degree two}
\begin{align*}
J&=p_1^2+2p_1p_2+\tfrac43p_2^2, &
Q&=\begin{bmatrix}1&1\\1&4/3\end{bmatrix}, & A&=\begin{bmatrix}1&1\end{bmatrix}, & b&=1.
\end{align*}
Substitute $p_1=1-p_2$: $J=1+\tfrac13p_2^2$.
Thus $\boldsymbol p^*=(1,0)^T$, $P^*(t)=t$, and $J^*=1$.
Allowing a quadratic does not lower the cost.
\subsection*{3. Degree three}
\[
Q=\begin{bmatrix}1&1&1\\1&4/3&3/2\\1&3/2&9/5\end{bmatrix},\qquad
A=\begin{bmatrix}1&1&1\end{bmatrix},\qquad b=1.
\]
The KKT equations $2Qp+A^T\lambda=0$, $Ap=b$ yield
$p^*=(1,0,0)^T$, $\lambda=-2$, so again $P^*(t)=t$, $J^*=1$.
This agrees with Euler--Lagrange: $P''=0$ for the integrand $(P')^2$.
In fact, Cauchy--Schwarz gives
\[
J\geq\left(\int_0^1 P'(t)dt\right)^2=(P(1)-P(0))^2=1.
\]
Equality requires constant velocity almost everywhere. Thus this answer is globally optimal for all polynomial degrees.
\subsection*{4. Add the terminal velocity $P'(1)=-2$}
The second constraint row is $[1,2,\ldots,N]$ and $b=(1,-2)^T$.
For a positive definite $Q$, the equality-constrained solution is
\[
p^*=Q^{-1}A^T(AQ^{-1}A^T)^{-1}b.
\]
\paragraph{(a) Degree two.}
\[
A=\begin{bmatrix}1&1\\1&2\end{bmatrix},\quad b=\begin{bmatrix}1\\-2\end{bmatrix},\quad
p^*=\begin{bmatrix}4\\-3\end{bmatrix}.
\]
Both coefficients are fixed by the constraints. Therefore
\[
\boxed{P^*(t)=4t-3t^2,\qquad J^*=\int_0^1(4-6t)^2dt=4.}
\]
\paragraph{(b) Degree three.}
\[
A=\begin{bmatrix}1&1&1\\1&2&3\end{bmatrix},\quad b=\begin{bmatrix}1\\-2\end{bmatrix}.
\]
Writing $p_3=z$ gives $p_1=4+z$ and $p_2=-3-2z$.
Substitution gives $J=4+z+\tfrac{2}{15}z^2$, minimized at $z=-15/4$.
Consequently
\[
\boxed{P^*(t)=\tfrac14t+\tfrac92t^2-\tfrac{15}{4}t^3,\qquad J^*=\tfrac{17}{8}.}
\]
Indeed $P(1)=1$ and $P'(1)=\tfrac14+9-\tfrac{45}{4}=-2$.
The cubic improves on the quadratic, but its cost remains above 1.
The unconstrained-degree infimum with this terminal velocity is 1 and is not attained:
velocity can change in an arbitrarily thin terminal layer. The Euler--Lagrange
linear solution cannot satisfy all three boundary conditions. The finite-degree QPs are well posed.

\section*{Deliverable 2: Multiple segments}
All counts below are per scalar coordinate, for fixed strictly increasing knot times
$t_0<\cdots<t_k$ and prescribed waypoint positions. Spatial coordinates can be solved independently.
\subsection*{1. Minimum snap ($r=4$)}
On each open segment, Euler--Lagrange gives $P_i^{(8)}=0$.
Each segment therefore has degree at most 7 and eight coefficients: $8k$ unknowns.
With position, velocity, acceleration and jerk fixed at the overall endpoints,
the conditions determining the stationary optimum are:
\begin{center}
\begin{tabular}{lll}\toprule
Type & Conditions & Scalar count\\\midrule
Waypoint & Both sides equal each interior position; two endpoints & $2k$\\
Free derivative & Match derivatives 1 through 6 at each interior knot & $6(k-1)$\\
Fixed derivative & Prescribe derivatives 1 through 3 at each endpoint & $6$\\\bottomrule
\end{tabular}
\end{center}
The total is $2k+6(k-1)+6=8k$. There are $k+1$ physical waypoint positions;
each interior waypoint supplies \emph{two} scalar equations.
For $k=3$: 6 waypoint equations, 12 free derivative equations and 6 fixed derivative equations give 24 equations.

Velocity, acceleration and jerk continuity are imposed for admissibility.
Continuity of snap, the fifth derivative and the sixth derivative follows from
stationarity when the common interior velocity, acceleration and jerk are free.
These last three are \emph{natural optimality conditions}, not additional
prescribed derivative values. The seventh derivative can jump because the
interior positions are fixed. In a QP one imposes only the waypoint, low-order
continuity and endpoint constraints; minimizing the cost supplies the natural conditions.

\subsection*{2. General derivative order $r\geq1$}
For $J=\sum_i\int_{t_{i-1}}^{t_i}(P_i^{(r)})^2dt$,
Euler--Lagrange yields $P_i^{(2r)}=0$. Hence there are $2r$ coefficients
per segment and $2rk$ unknowns.
\begin{itemize}
\item Waypoint equations: $P_1(t_0)=\bar p_0$, $P_k(t_k)=\bar p_k$,
and $P_i(t_i)=P_{i+1}(t_i)=\bar p_i$ at the $k-1$ interior knots: $2k$ equations.
\item Fixed endpoint derivatives: orders $1,\ldots,r-1$ at each endpoint:
$2(r-1)$ equations (typically all zero).
\item Free derivative equations: continuity of orders $1,\ldots,2r-2$ at
interior knots: $(2r-2)(k-1)$ equations. Orders $1,\ldots,r-1$ are
admissibility constraints; orders $r,\ldots,2r-2$ are natural optimality conditions.
\end{itemize}
Their total is
\[
\boxed{2k+(2r-2)(k-1)+2(r-1)=2rk.}
\]
To see the natural conditions, integrate the first variation by parts $r$ times:
\[
\tfrac12\delta J_i=
\left[\sum_{j=0}^{r-1}(-1)^jP_i^{(r+j)}\delta P_i^{(r-1-j)}\right]_{t_{i-1}}^{t_i}
+(-1)^r\int_{t_{i-1}}^{t_i}P_i^{(2r)}\delta P_i\,dt.
\]
At an interior knot $\delta P=0$ since position is fixed. The common variations
$\delta P^{(1)},\ldots,\delta P^{(r-1)}$ are arbitrary, giving continuity of orders
$r,\ldots,2r-2$. The derivative of order $2r-1$ need not be continuous.

\paragraph{Relation to the code.}
The supplied library uses $N=10$ coefficients (degree at most 9) and
\texttt{makeStartOrEnd(..., SNAP)} fixes derivatives through order 4 to zero.
This is a richer polynomial parameterization with an extra endpoint snap condition;
it differs from the minimal degree-7, endpoint-through-jerk theoretical count above.
The implementation follows the supplied library convention.
\paragraph{Source.}
Problem statement: \url{https://vnav.mit.edu/labs_2023/lab4/exercises.html}.
All calculations and explanations above were derived for these exercises.
\end{document}
